\section{Attention 机制：注意力是你所需要的一切}

\subsection{超越循环：寻找新范式}

\begin{figure}[H]
\centering
\includegraphics[width=0.85\textwidth]{slides-images/slide-60.jpg}
\caption{序列建模的目标：从输入序列到特征再到输出\protect\footnotemark}
\end{figure}
\footnotetext{Slides 第60页。}

让我们回到序列建模的根本目标：给定输入序列，提取特征，生成预测。RNN 通过逐步处理来捕捉序列依赖，但如果我们能\textbf{同时看到整个序列}，并自动识别其中重要的部分，是否能做得更好？

一种朴素的想法是：将所有输入拼接成一个大向量，用前馈网络处理。但这有两个问题：
\begin{itemize}
    \item 参数规模会随序列长度爆炸式增长，不具有可扩展性
    \item 拼接操作破坏了位置/顺序信息
\end{itemize}

\begin{figure}[H]
\centering
\includegraphics[width=0.75\textwidth]{slides-images/slide-62.jpg}
\caption{朴素方法的问题：拼接输入丧失位置信息\protect\footnotemark}
\end{figure}
\footnotetext{Slides 第62页。}

\subsection{什么是 Attention}

\textbf{Attention}（注意力）机制的核心思想源自人类认知：当我们观察一个场景时，不会均匀地关注所有内容，而是自动聚焦于最相关、最重要的部分。

\begin{importantbox}{Attention 的核心定义}
Attention 是一种让网络能够\textbf{自动识别并聚焦于输入中最重要部分}的机制。它通过计算输入不同部分之间的\textbf{相关性得分}（Attention Score），来决定在生成输出时应该``关注''哪些输入。
\end{importantbox}

\subsection{搜索引擎类比：Query-Key-Value}

讲者用一个非常直观的搜索引擎类比来解释 Attention 的工作原理：

\begin{figure}[H]
\centering
\includegraphics[width=0.75\textwidth]{slides-images/slide-67.jpg}
\caption{搜索引擎类比：Query-Key-Value 框架\protect\footnotemark}
\end{figure}
\footnotetext{Slides 第67页。}

假设你在视频搜索引擎中搜索``deep learning course''（Query）：
\begin{enumerate}
    \item 数据库中每个视频都有一个标签/描述（Key）
    \item 搜索引擎计算 Query 与每个 Key 的\textbf{相似度}
    \item 找到最匹配的 Key 后，提取对应的视频内容（Value）
\end{enumerate}

\begin{knowledgebox}{Query-Key-Value 的三元组}
Attention 机制中的三个核心概念：
\begin{itemize}
    \item \textbf{Query (Q)}：你正在寻找什么——当前需要被回答的``问题''
    \item \textbf{Key (K)}：数据库中每个条目的``标签''——用于匹配的索引
    \item \textbf{Value (V)}：与每个 Key 关联的实际内容——匹配成功后要提取的信息
\end{itemize}
Attention 的过程就是：用 Query 在 Keys 中搜索最匹配的条目，然后提取对应的 Values。
\end{knowledgebox}

\subsection{本章小结}

Attention 机制的核心是一个``搜索与检索''的过程：通过计算 Query 和 Key 之间的相似度来确定应该关注哪些输入，然后从 Value 中提取相关信息。这种机制不需要逐步处理序列，可以直接建模任意两个位置之间的依赖关系。

\newpage
